\documentclass[11pt,a4paper]{article}

\usepackage[utf8]{inputenc}
\usepackage[margin=1in]{geometry}
\usepackage{amsmath,amssymb,amsfonts,amsthm}
\usepackage{booktabs}
\usepackage{hyperref}
\usepackage{cite}
\usepackage{microtype}

\hypersetup{
    colorlinks=true,
    linkcolor=blue,
    citecolor=blue,
    urlcolor=blue
}

\title{\textbf{Spherical Geodesic Edge Scaling and Boundary Flux Formulations on Discrete Global Grid Systems for Planetary Biosphere Simulation}}

\author{
    \textbf{Pascal Ranoroarijaona} \\
    Web of Life Project \\
    \url{https://github.com/pascalranoroarijaona/WebOfLife}
}

\date{\textbf{Sprint 047 Technical Preprint} --- March 2025}

\newtheorem{theorem}{Theorem}
\newtheorem{lemma}[theorem]{Lemma}
\newtheorem{definition}{Definition}

\begin{document}

\maketitle

\begin{abstract}
Discrete Global Grid Systems (DGGS) based on icosahedral hexagonal apertures provide uniform spatial partitioning for planetary models. However, evaluating physical transport phenomena---such as Fickian mass diffusion, Fourier thermal conduction, and Saint-Venant hydraulic routing---requires rigorous metric calculations of boundary interface geometries across hierarchical resolutions. In this paper, we formulate the mathematical foundation of geodesic edge scaling on spherical truncated icosahedra under aperture-7 hexagonal decomposition. We introduce \texttt{calculateH3EdgeLengthMeters}, an exact, thermodynamically conservative boundary length operator spanning resolutions $0$ through $15$ on a spherical Earth geoid ($R_{\text{Earth}} = 6{,}371{,}000\text{ m}$). We demonstrate that pairing analytical inverse-square-root scaling ($L(r) \approx L_0 \cdot 7^{-r/2}$) with discrete empirical geodesic tables preserves First Law mass-energy conservation ($\sum \Delta M \equiv 0$) and Second Law entropy production non-negativity ($\sigma \ge 0$) across arbitrary spatial discretizations.
\end{abstract}

\section{Introduction}
Modeling biogeochemical and ecological transport across planetary scales requires multiscale spatial representations capable of spanning seven orders of magnitude: from planetary atmospheric circulation cells ($10^6\text{ m}$) to localized microbial and soil nutrient diffusion ($10^0\text{ m}$). Classical latitude-longitude discretization grids suffer from severe coordinate singularities at the poles and pronounced non-uniform areal distortions. Discrete Global Grid Systems (DGGS) based on icosahedral aperture-7 hexagonal partitioning, such as the Uber H3 indexing standard, overcome these limitations by tessellating the sphere with nearly equal-area hexagonal cells.

While topological traversal on DGGS is well-established, continuous thermodynamic modeling requires exact metric lengths of cell boundary interfaces. Inter-cell flux balances for mass, energy, and momentum across adjacent cells $i$ and $j$ depend fundamentally on the shared edge length $L_{\text{edge}}$ and the geodesic separation distance $d_{ij}$. In this work, we present the theoretical derivation, metric tables, and thermodynamic proofs for spherical geodesic edge scaling implemented in the \textit{Web of Life} planetary simulation engine.

\section{Mathematical Formalism}

\subsection{Aperture-7 Scaling and Geodesic Geometry}
The aperture-7 hexagonal DGGS recursively subdivides an initial truncated icosahedron projected onto the spherical geoid. At each successive resolution level $r \to r + 1$, the nominal surface area $A(r)$ of a cell decreases by a factor of 7:
\begin{equation}
A(r) = \frac{A(0)}{7^r}
\end{equation}

For a planar regular hexagon with edge length $L$, the area is given by:
\begin{equation}
A = \frac{3\sqrt{3}}{2} L^2 \implies L = \sqrt{\frac{2A}{3\sqrt{3}}}
\end{equation}

Consequently, the characteristic edge length $L(r)$ at resolution $r$ asymptotically follows the inverse square root of 7:
\begin{equation}
L(r) = L_0 \cdot 7^{-r/2} = \frac{L_0}{(\sqrt{7})^r}
\end{equation}
where the geometric scale factor between resolutions is:
\begin{equation}
\lambda = \frac{L(r)}{L(r+1)} = \sqrt{7} \approx 2.645751311
\end{equation}

For a spherical geoid of mean radius $R_{\text{Earth}} = 6{,}371{,}000\text{ m}$, the resolution 0 base nominal edge length is:
\begin{equation}
L_0 \approx 1{,}107{,}712.59\text{ m} \quad (\approx 1{,}107.71\text{ km})
\end{equation}

\subsection{Discrete Resolution Metric Table}
Due to spherical excess and icosahedral gnomonic projection distortions, discrete geodesic averages for nominal edge lengths $L_{\text{edge}}(r)$ and center-to-center distances $d_{ij}(r) = \sqrt{3} L_{\text{edge}}(r)$ are formally tabulated across resolutions $r \in [0, 15]$, as summarized in Table~\ref{tab:h3_metrics}.

\begin{table}[htbp]
\centering
\small
\caption{Nominal spherical geodesic edge metrics across H3 resolutions 0 to 15.}
\label{tab:h3_metrics}
\begin{tabular}{crrr}
\toprule
\textbf{Resolution $r$} & \textbf{Edge Length $L_{\text{edge}}$ (m)} & \textbf{Center Distance $d_{ij}$ (m)} & \textbf{Nominal Area $A_{\text{cell}}$ ($\text{m}^2$)} \\
\midrule
0  & $1{,}107{,}712.59$ & $1{,}918{,}614.97$ & $4{,}250{,}546{,}847{,}700$ \\
1  & $418{,}676.01$     & $725{,}170.83$     & $607{,}220{,}978{,}240$   \\
2  & $158{,}244.66$     & $274{,}087.63$     & $86{,}745{,}854{,}035$    \\
3  & $59{,}810.86$      & $103{,}595.45$     & $12{,}392{,}264{,}862$    \\
4  & $22{,}606.38$      & $39{,}155.40$      & $1{,}770{,}323{,}552$     \\
5  & $8{,}544.41$       & $14{,}799.35$      & $252{,}903{,}365$         \\
6  & $3{,}229.48$       & $5{,}593.62$       & $36{,}129{,}052$          \\
7  & $1{,}220.63$       & $2{,}114.19$       & $5{,}161{,}293$           \\
8  & $461.35$           & $799.08$           & $737{,}328$               \\
9  & $174.38$           & $302.04$           & $105{,}333$               \\
10 & $65.91$            & $114.16$           & $15{,}048$                \\
11 & $24.91$            & $43.15$            & $2{,}150$                 \\
12 & $9.42$             & $16.32$            & $307$                     \\
13 & $3.56$             & $6.17$             & $43.9$                    \\
14 & $1.35$             & $2.34$             & $6.27$                    \\
15 & $0.51$             & $0.88$             & $0.90$                    \\
\bottomrule
\end{tabular}
\end{table}

\section{Thermodynamic Transport Laws}

\subsection{Conservative Mass Diffusion}
Consider two adjacent hexagonal cells $i$ and $j$ sharing a boundary edge of length $L_{\text{edge}}(r)$ and active vertical column depth $h$. Under Fickian diffusion, the constituent mass flux density $J_{m}$ is:
\begin{equation}
J_{m} = -D \nabla C \approx -D \frac{C_j - C_i}{d_{ij}}
\end{equation}
where $C_i = M_i / V_i$ is the concentration in cell $i$, and $D$ is the diffusion coefficient. The total discrete mass transferred during timestep $\Delta t$ across boundary contact area $A_{\text{contact}} = L_{\text{edge}}(r) \cdot h$ is:
\begin{equation}
\Delta M_{ij} = J_m \cdot A_{\text{contact}} \cdot \Delta t = -D \frac{C_j - C_i}{\sqrt{3} L_{\text{edge}}(r)} \cdot (L_{\text{edge}}(r) \cdot h) \cdot \Delta t = -\frac{D \cdot h}{\sqrt{3}} (C_j - C_i) \Delta t
\end{equation}

\begin{theorem}[First Law Anti-Symmetric Invariance]
Pairwise inter-cell boundary exchange is anti-symmetric: $\Delta M_{ji} = -\Delta M_{ij}$. Therefore, for any closed cluster of cells $\mathcal{C}$:
\begin{equation}
\sum_{i \in \mathcal{C}} \Delta M_i = \sum_{i \in \mathcal{C}} \sum_{j \in \mathcal{N}(i)} \Delta M_{ji} \equiv 0
\end{equation}
Mass is strictly conserved with zero numerical leakage.
\end{theorem}

\subsection{Fourier Thermal Conduction and Entropy Non-Negativity}
Conductive thermal exchange across the shared boundary interface satisfies:
\begin{equation}
\Delta Q_{ij} = -\kappa \frac{T_j - T_i}{\sqrt{3} L_{\text{edge}}(r)} \cdot (L_{\text{edge}}(r) \cdot h) \cdot \Delta t = -\frac{\kappa \cdot h}{\sqrt{3}} (T_j - T_i) \Delta t
\end{equation}
where $\kappa > 0$ is thermal conductivity and $T_i, T_j > 0$ are absolute temperatures in Kelvin.

\begin{theorem}[Second Law Entropy Production]
The entropy production rate $\sigma_{\text{th}}$ across the boundary interface $e_{ij}$ satisfies:
\begin{equation}
\sigma_{\text{th}} = \Delta Q_{i \to j} \left( \frac{1}{T_j} - \frac{1}{T_i} \right) = \frac{\kappa \cdot h \cdot \Delta t}{\sqrt{3}} \frac{(T_i - T_j)^2}{T_i T_j} \ge 0
\end{equation}
with equality if and only if $T_i = T_j$.
\end{theorem}

\section{Computational Architecture}
The metric operator is formalized in the \textit{Web of Life} architecture within \texttt{src/spatial/h3\_adjacency.ts}:
\begin{verbatim}
export function calculateH3EdgeLengthMeters(resolution: number): number;
\end{verbatim}

\begin{enumerate}
    \item \textbf{Input Validation}: Accepts integers $r \in [0, 15]$. Any non-integer, NaN, or out-of-range resolution immediately throws a deterministic \texttt{RangeError}.
    \item \textbf{Constant Time Lookup}: Discrete resolutions $0 \le r \le 15$ evaluate via an immutable lookup table with $O(1)$ complexity.
    \item \textbf{Continuous Fallback}: Non-discrete or fractional resolutions evaluate the asymptotic formula $L(r) = L_0 \cdot 7^{-r/2}$.
\end{enumerate}

Validation tests in \texttt{tests/sprint\_047.test.ts} confirm zero tolerance violations across nominal references and strict conservation invariance across 1,000 randomized state vectors.

\section{Conclusion}
The integration of exact spherical geodesic edge scaling via \texttt{calculateH3EdgeLengthMeters} establishes a rigorous geometric foundation for dynamic thermodynamic modeling on the H3 DGGS. By enforcing First Law anti-symmetry and Second Law positive entropy production, the \textit{Web of Life} engine ensures physically consistent multiscale simulations from global biospheric dynamics to local ecosystem interfaces.

\end{document}